\section{Clasificación completa de partículas y estados físicos}
\label{sec:catalogo-realizaciones-tipadas-rev8}
\index{sectores físicos!inventario tipado}
\index{atlas!81/56/13/324}

Las realizaciones se ordenan por objeto, multiplicidad y codominio. El dominio matemático contiene \(81\) posiciones, agrupadas en
codominios distintos. El dominio matemático contiene \(81\) posiciones, agrupadas en
\(56\) familias de trayectoria y \(13\) clases de extensión; las \(324\)
trayectorias orientadas añaden la orientación inicial y no representan masas
ni partículas adicionales. La columna ``operador de evaluación'' no impone una masa donde el objeto exige
una carta nula, una composición previa, un invariante topológico o una
obstrucción de prolongación.

\subsection{Definiciones espectrales de partícula, espín, color y sabor}

Una \emph{partícula} es una sección superviviente realizada como ruta
orientada, dotada de registro, representación de simetría, operador espectral
y carta física. Un fermión realiza la graduación exterior y las relaciones
CAR; un bosón realiza la completación simétrica y las relaciones CCR. El
espín registra la holonomía de la orientación: el retorno de \(2\pi\) cierra
la fase bosónica, mientras la sección espinorial recupera su estado tras
\(4\pi\). Esta distinción enlaza la doble hoja areal--volumétrica con la
paridad y la exclusión de Pauli.

El \emph{color} es la representación ternaria que eleva la coordenada
\(r-c\pmod3\) a la órbita fundamental y, después, al adjunto de
\(\mathfrak{su}_3\). El \emph{sabor} identifica una base de realización
dentro de una representación con carga, color y espín fijados; CKM y PMNS
transportan entre la base de interacción y la base espectral. Los neutrinos
son secciones fermiónicas neutras: sus etiquetas de sabor pertenecen al
lector de mezcla y sus masas a los autovectores del operador correspondiente.

Una partícula virtual es una sección finita con registro y propagador hasta
un horizonte de prolongación; una realización Majorana es el sector
fermiónico real fijado por conjugación; y los sectores de Fibonacci, Ising y
\(\mathbb Z/6\mathbb Z\) publican dimensiones cuánticas, reglas de fusión y
caracteres de trenza. Cada definición conserva así el codominio que la
estructura le asigna.

